% \input MET WHILE A FILE NAME IS BEING SCANNED DOES NOT OPEN ITS FILE.
%
% It stands a \relax in front of itself and waits. The \relax ends
% whatever was being scanned, and the \input is met again afterwards:
%
%   \font\zz=nosuchfont\input tripfile
%
%   {\font}
%   {\input}                     <- expanded here, and it waits
%   ! Font \zz=nosuchfont not loadable: Metric (TFM) file not found.
%   <to be read again> \relax
%   {\relax}                     <- the \relax it stood in front of itself
%   {\input}                     <- met again, and NOW the file opens
%   (./tripfile.tex
%
% The name being scanned is not only the font's own: the size that may
% follow it counts too, and that is where this \input is met -- the font
% name ended at the control sequence, and the \input turned up while
% `at' was being looked for.
%
% \input IS EXPANDED, NOT OBEYED, which is what puts its trace line above
% the fault rather than below it. hstex had it as a command of the main
% loop, so it could only be traced after the font had been looked for and
% failed.
%
% A CHARACTER THE FONT DOES NOT HAVE DOES NOT CARRY A RUN. A run of
% characters is traced once, at its first -- but a character that joins no
% list carries nothing, so the next one is traced too. With the font above
% left at \nullfont, `hello' draws five lines where a real font draws one.
%
% RUN THIS UNDER `pdftex -ini'; it sets its own catcodes and wants a file
% called tripfile.tex beside it.
\catcode`\{=1 \catcode`\}=2 \catcode`\#=6
\tracingonline=1 \tracingcommands=2
\message{[A a font that will not load, then an input]}
\font\zz=nosuchfont\input tripfile
\message{[B]}
\tracingcommands=0
\end
